\section*{3月21日：勾股定理简单应用}


\begin{enumerate}
\item 解答：

（1）在直角三角形中，两条直角边分别为 $3$ 和 $4$，求斜边的长度。

（2）直角三角形的一条直角边为 $5$，斜边为 $13$，求另一条直角边的长度。

（3）直角三角形斜边长为 $10$，一条直角边长为 $6$，求另一条直角边长。

（4）一个直角三角形的两条直角边分别为 $9$ 和 $12$，求斜边长。

\item 解答：

（1）已知一个直角三角形的两边长分别为 $8$ 和 $10$，求第三边的长度。

（2）一个直角三角形的一边长为 $7$，另一边长为 $25$，求第三边的长。
\end{enumerate}






\section*{3月23日：勾股定理与立体图形最短路径}

\begin{enumerate}
\item 小红将一条彩带缠绕圆柱，正好从$A$点绕到正上方的$B$点
，已知易拉罐的底面周长为$12cm$，高为$16cm$，则所需彩带最短为$\_\_\_\_\_\_.$
\item 已知圆柱底面的周长为$6dm$，圆柱高为$4dm$，在圆柱的侧面有一只蚂蚁，从$A$到$B$，再从$B$回到$A$，恰好爬行一圈，则这只蚂蚁爬行的最短路程为$\_\_\_\_\_\_.$
\end{enumerate}


\newpage

\section*{3月24日：勾股定理花式求最值}